\section{Vote}
\label{vote}

\begin{enumerate}
 \item Le vote aura lieu au scrutin secret en utilisant un système de vote transférable.
  \begin{enumerate}
   \item Pour les élections à un poste unique, la méthode du vote à second tour instantané (VSTI) sera utilisée.
    \begin{enumerate}
     \item Dans l'éventualité où aucune personne candidate n'obtient plus de 50 \% des votes suite à cette procédure, aucune personne candidate ne sera déclarée élue ou nominée, et une nouvelle élection pourra être tenue.
    \end{enumerate}
   \item Lorsque deux postes ou plus doivent être pourvus, les postes seront pourvus séquentiellement en utilisant la méthode du VSTI. Suivant l'élection ou la nomination de chaque personne candidate, cette personne candidate sera retirée des bulletins de vote restants, et le processus sera répété jusqu'à ce que tous les postes aient été pourvus.
    \begin{enumerate}
     \item Une personne candidate doit obtenir plus de 50 \% des votes pour être déclarée élue ou nominée, selon le cas.
    \end{enumerate}
   \item L'exécutif national soumettra une liste classée des personnes candidates à la DGE pour être utilisée uniquement en cas d'égalité.
  \end{enumerate}
 \item Les procédures d'administration et de dépouillement des votes seront établies par la DGE conformément à ces statuts.
 \item Les votes seront comptés par la DGE.
  \begin{enumerate}
   \item L'assemblée générale peut nommer, par résolution, une personne scrutatrice, qui supervisera le dépouillement.
  \end{enumerate}
\end{enumerate}
